\chapter{Alineaci\'on con Green Tech}

EcoBrief Bolivia encaja en Green Tech porque promueve un uso m\'as eficiente y responsable de software, datos e inteligencia artificial.

\section{Reducci\'on de navegaci\'on innecesaria}

El usuario no necesita abrir m\'ultiples p\'aginas para entender los hechos principales. El sistema entrega una s\'intesis priorizada que reemplaza la navegaci\'on dispersa. Cada resumen le\'ido representa p\'aginas enteras que no fueron cargadas.

\section{Reducci\'on de duplicaci\'on informativa}

La deduplicaci\'on evita que la misma historia sea procesada y mostrada varias veces. Esto reduce el almacenamiento, el ancho de banda y las llamadas a la IA para contenido redundante.

\section{Reducci\'on de llamadas IA redundantes}

La IA se usa despu\'es de limpiar, filtrar, deduplicar y rankear. Esto evita enviar contenido irrelevante o repetido al modelo. Se estima una reducci\'on de hasta el 65\,\% en el volumen de contenido procesado por IA.

\section{Reducci\'on de consumo de datos}

El usuario puede leer res\'umenes compactos en lugar de cargar p\'aginas completas con im\'agenes, publicidad, \textit{trackers} y scripts. Se estima un ahorro de aproximadamente 0.8 MB por p\'agina evitada.

\section{Promoci\'on del consumo digital responsable}

Las preferencias por categor\'ia y frecuencia reducen la informaci\'on no solicitada y evitan el \textit{spam} informativo. El usuario elige qu\'e recibir y cu\'ando.

\section{Reducci\'on de dependencia del scroll social}

EcoBrief ofrece una alternativa al consumo informativo en redes sociales. En lugar de depender de TikTok, Facebook u otros \textit{feeds} con videos y contenido no siempre trazable, el usuario recibe res\'umenes compactos, priorizados y enlazados a fuentes identificadas.

\section{Trazabilidad y confianza informativa}

El sistema conserva el enlace a la fuente original de cada noticia. Prioriza art\'iculos con cuerpo, fecha y fuente identificada. Reduce la exposici\'on a contenido no trazable que circula en redes sociales sin contexto ni verificaci\'on. No reemplaza el \textit{fact-checking} period\'istico, pero ayuda a evitar la dependencia de informaci\'on viral sin fuente clara.
